\documentclass[10pt]{article}
\usepackage[a4paper,margin=25mm]{geometry}
\usepackage{amsmath,amssymb,amsthm,mathrsfs,longtable,hyperref}
\emergencystretch=3em
\newtheorem{theorem}{Theorem}[section]
\newtheorem{lemma}[theorem]{Lemma}
\newtheorem{proposition}[theorem]{Proposition}
\newtheorem{corollary}[theorem]{Corollary}
\newcommand{\PSH}{\operatorname{PSH}}
\newcommand{\FS}{\operatorname{FS}}
\newcommand{\vol}{\operatorname{vol}}
\newcommand{\MA}{\operatorname{MA}}
\newcommand{\E}{\mathcal E}
\newcommand{\OO}{\mathcal O}
\newcommand{\dd}{dd^c}
\title{Bounded-obstacle section approximation and ordered energy-gap tails\\
\large Mathematical progress, with exact conditions, for original P16.5}
\author{An unverified research reconstruction from the two caller grants}
\date{6 October 2026}
\begin{document}
\maketitle
\begin{abstract}
We reconstruct genuine closure, continuous-obstacle envelopes, and the
section-volume energy of the single-net numerical cone from the two supplied
grants. An actual ample-gap tensor argument works for every bounded continuous
obstacle, not only the zero obstacle. Both full-section envelope discrepancies
therefore vanish outside one fixed actual Kodaira divisor. We verify the
small-cost capped-section construction of a genuine finite-energy pole and
algebraic nullness of every intrinsic target. We also prove a finite-nef
ordered energy-gap estimate. Under an explicitly stated, still unaudited
outer-output and full-moving compatibility package, it gives uniform
approximation tails on energy-bounded output families, absolute original
continuous $L^1$ approximants, relative weighted continuity, and whole
supporting planes. This conditional output theorem does not imply strong
integration against every intrinsic target. Arbitrary-target attainment,
strict Borel comparison, strict separation, and uniqueness remain unproved.
The complete original P16.5 is not solved; the canonical result is
\texttt{candidate=false}. The additional zero-volume equation theorem is
exactly the numerical-zero class.
\end{abstract}

\section{Exact domain, premises, and boundary of the claims}
Let $k$ be any algebraically closed trivially valued field of characteristic
zero, $X$ any smooth integral projective $k$-variety of dimension $n\geq1$,
$K=X^{\rm an}$, and $L\in\operatorname{Pic}(X)\otimes\mathbb Q$ any actual
big rational bundle. Put $\theta=c_1(L)$ and $V=\vol_X(L)>0$.
There is no countability, toric, skeletal, support, or section-ring
finite-generation hypothesis. Original supported semivaluations remain
points of $K$. Actual tensor identities are distinguished from numerical
intersection comparisons. Every rational tensor power used for sections
is cleared. Scalars use the canonical trivial metric on the indicated
bundle; $v_0$ is the generic trivial valuation.

G0 is the entire supplied ample P16.1 theorem, only in its actual ample
rational-bundle domain and exact retained hypotheses. G1 is the entire
supplied resolved relative-generation theorem and its one generically
finite FS net, only in its specified signature. No nonample equation,
unbounded estimate, or general closure is granted. The supplied Full151,
Full154, Root195, Root196, Root197 and their ledgers are unreviewed
mathematical data. This paper does not use their asserted conclusions
without a proof here or an explicit conditional hypothesis below.

Ordinary background is characteristic-zero equivariant resolution and
principalization, model realization and density of rational PL functions,
Kodaira decomposition, eventual ample generation and openness of the ample
cone, normal proper birational section descent, section-growth limits, and
the algebraic positive products and volume differential of BFJ. These are
ordinary algebraic inputs, not nonample NA equation grants. The precise
primary units and applications are recorded in the appendix. All new
measure-theoretic passage arguments used below are proved here.

\paragraph{What is proved and what is conditional.}
Sections 2--6 give proof drafts, relative to G0, G1 and the scoped ordinary
background, for the actual cone, envelopes, scalar energy, and target
algebraic-nullness statements. Section 7 is finite-model algebra and Hodge
calculus, with the corrected coefficient orientation. Sections 8--10 are
conditional on package F, stated in Section 8. In particular we have not
independently completed the entire construction of an outer-output map
from Full151, and do not identify that proposed map as the entire genuine
finite-energy operator required by the task. Stating F makes this
limitation an antecedent of every output theorem, not a hidden grant.
Section 11 gives precise remaining equation assertions. Neither compiled
PDFs nor arithmetic checks certify any of these proof drafts.

\section{The actual single-net cone and its elementary operations}
A full FS scalar on an actual rational $B$ is
\[
 h=\frac1m\max_j\{\log\|s_j\|_{\rm triv}+\lambda_j\},\qquad
 0\ne s_j\in H^0(X,mB),\quad\lambda_j\in\mathbb R.
\]
Common zeros are permitted; these functions are extended-continuous and
generically finite. Define $\PSH(\theta)$ as the all-point decreasing
limits of one net of such functions on actual $B_i$ with
$c_1(B_i)\to\theta$, whose generic limiting value is finite. Thus every
element is an actual usc function on $K$, rather than a completion point.
This definition does not itself grant twice-decreasing closure.

For an actual generated bundle a weight-zero generating system has scalar
zero everywhere: at the valuation-ring center one generating section is a
unit, including at supported points. Tensoring such a system changes the
actual bundle but preserves the scalar. If $\Omega-L$ is ample, a tail of
any defining net has $\Omega-B_i$ ample. The scalars embed into the one
actual ample cone of $\Omega$. To make its stages finite append an ample
generating system at weight $-N$ and use the product net
$\max(h_i,-N)$. G0 supplies genuine ample membership, finite divisorial
values, and all-point divisorial order determination for the endpoint.
This asserts no equality of ample and big energies.

Canonical section norms are at most one and are one at $v_0$ for nonzero
sections. The supremum of each stage and its endpoint is its generic value.
For two defining nets, choose valid triples $(i,j,\epsilon)$ with
$L+\epsilon A-B_i$ and $L+\epsilon A-C_j$ ample, where $A$ is fixed ample.
The valid triples are directed by taking a smaller rational $\epsilon>0$
and sufficiently late indices. Tensoring into the same actual target and
uniting sections gives a decreasing FS net for the maximum, with classes
tending to $\theta$. Constants and rational scaling preserve single nets.

Convexity with real coefficients is also direct. Put both input nets below
one finite $C$. Choose positive rational $r_j\uparrow r$ and
$s_j\uparrow1-r$ for $0<r<1$. The FS scalars
\[
 r_j(h_i-C)+s_j(g_l-C)+C
\]
decrease on the product index and have the desired convex limit. Their
actual bundle is $r_jB_i+s_jC_l$ and their numerical class tends to
$\theta$. No subtraction of poles occurs. Endpoint coefficients zero are
handled separately. Integer multiplication follows from tensoring section
systems. These operations require no general decreasing closure.

\section{Every bounded obstacle has exact section minorants off a divisor}
Fix once an actual Kodaira identity
\[
 L=aA+E_0,\qquad a>0\text{ rational},\quad A\text{ ample},\quad E_0\geq0.
\]
Set $W=q_{E_0}=\log\|s_{E_0}\|_{\rm triv}\leq0$ and
$Z=E_0^{\rm an}$. $W$ is finite continuous exactly off $Z$ and is a
single-section scalar on the actual $L-aA$. Tensoring with a generated
power of $aA$ also realizes its scalar on $L$.
For bounded continuous $f$ put
\[
 Q_L(f)=\sup\{h\in\FS(L):h\leq f\},\qquad
 P_\theta(f)=\sup\{u\in\PSH(\theta):u\leq f\}.
\]
For now the second expression is a raw supremum, not assumed genuine.
The minorant $W-\|f\|_\infty$ gives a finite lower bound off $Z$.

\begin{lemma}[Actual ample-gap approximation]\label{sections}
For every genuine $u\leq f$, every $x\notin Z$ satisfies
$u(x)\leq Q_L(f)(x)$. Consequently $P_\theta(f)=Q_L(f)$ off $Z$.
No envelope or decreasing-closure theorem is needed for this lemma.
\end{lemma}
\begin{proof}
Take its literal defining FS net $h_i\downarrow u$ on actual $M_i$.
Fix rational $0<t<1$ and $\epsilon>0$ before taking a tail. The actual
bundle
\[
 C_{i,t}=L-(1-t)M_i-t(L-aA)
\]
has numerical class tending to $ta[A]$, so is ample for a tail. Its
generated-system scalar is zero. Since $f$ is continuous, the compact
closed sets $\{h_i\geq f+\epsilon\}$ decrease to the empty set. A tail
therefore satisfies $h_i<f+\epsilon$ everywhere. Tensor the literal
sections, clearing rational $t$ and every bundle denominator, to obtain
an FS scalar on the actual $L$:
\[
 h_{i,t,\epsilon}=(1-t)h_i+tW-\epsilon-t\|f\|_\infty.
\]
Its difference from $f$ is at most
$-t f+tW-t\|f\|_\infty-t\epsilon\leq0$; the inequality holds at
all supported poles too. Thus $Q_L(f)\geq h_{i,t,\epsilon}$.
At $x\notin Z$ let $i$ tend to its defining endpoint, then
$\epsilon\downarrow0$, then rational $t\downarrow0$. If $u(x)=-\infty$
the desired inequality is immediate; otherwise all expressions here are
finite. The reverse envelope inequality follows from actual FS membership.
\end{proof}

The normalization of incoming stages is unnecessary. In particular a
constant shift of the original obstacle remains part of the minorant.
This is the extension of Root196's zero-obstacle argument to every
bounded continuous obstacle. It also reconstructs the pertinent conclusion
of the primary Addendum, Lemma 1.5, from the present single-net definition
rather than assuming its nonample Theorem 4.15 application.

For a rational model obstacle $f_D$, G1 supplies a single genuine endpoint
$\psi$ with $Q_L(f_D)\leq\psi\leq f_D$. Any genuine competitor $u\leq f_D$
is at most $Q_L(f_D)$ off $Z$ by Lemma \ref{sections}. Hence $u\leq\psi$
on every rational divisorial valuation. Embed both into one actual ample
cone and apply G0's all-point order determination. Thus
\[
 P_\theta(f_D)=\psi\in\PSH(\theta).
\tag{3.1}
\]
This uses only G1's one net, not nonample decreasing closure. In particular
$R=P_\theta(0)$ is genuine, $R(v_0)=0$, and every normalized genuine $u$
satisfies $u\leq R\leq0$. For any big actual $B$ the same construction
gives $R_B$. A late G1 reference stage, normalized as $\psi_B\leq0$ on
an actual $B+C/m$, obeys
\[
 R_B\leq\psi_B+c_B
\tag{3.2}
\]
for its finite generic shift $c_B$. It is not claimed that the normalized
stages themselves decrease or lie above $R_B$ without that shift.

\section{Reconstruction of genuine closure and continuous envelopes}
We include the flattening argument because G1 alone does not grant it.
Suppose $F$ is FS and $F\geq(1-\delta)\psi-k$, where $\psi\leq0$ is FS
and $0<\delta<1$. If $h$ is usc and $h\leq\psi+C$, extend $h-F$ by
$-\infty$ at poles of $F$. This is an usc comparison extension:
off those poles it is at most $\delta\psi+C+k$, which tends to
$-\infty$ at them. A pole of $F$ forces $\psi=h=-\infty$.
This assigned value is not arithmetic subtraction of infinities.
For a decreasing net with common $C$, empty intersection of its
nonnegative comparison superlevels gives an entirely negative tail by
compactness.

Let $u_i\downarrow u$ in $\PSH(\theta)$, with $u(v_0)$ finite. Restrict
to a tail $u_i\leq M$. Fix actual ample $A$. For rational $\eta>0$ set
$T_\eta=L+\eta A$. Choose $0<\epsilon<\eta/4$ and a normalized G1
reference stage $\psi$ on an actual $B+C/m$ with $B=L+\epsilon A$, late
enough that $T_\eta-(B+C/m)$ is ample. Choose small rational $\delta>0$
so $T_\eta-(1-\delta)(B+C/m)$ remains ample. A special stage is an FS
scalar $F$ in the specified actual $T_\eta$ such that
\[
 F\geq u+\kappa,\qquad F\geq(1-\delta)\psi-k,
\tag{4.1}
\]
with rational $\kappa>0$ and finite $k$. It exists: tensor a late defining
majorant of some $u_i$, and the floor, into $T_\eta$ and take their maximum.

For each fixed old reference $B_l,\psi_l$, every $u_i$ and sufficiently
late defining stages $h$ of any fixed $u_i$ embed into $B_l$. If the
stage generic values are at most $M+1$, minimality and (3.2) give the
common bound
\[
 h\leq R_{B_l}+M+1\leq\psi_l+c_l+M+1.
\tag{4.2}
\]
This is the extra old-pole bound absent from an abstract twice-closure
argument.

Given finitely many old special stages $F_l$ and guards $\kappa_l$, choose
new $\eta$ arbitrarily small below all old $\epsilon_l$. Choose its new
reference stage $\psi$ on $S$ sufficiently close numerically to $L$ that
$B_l-S$ is ample for every $l$. Normalization gives
$\psi\leq R_{B_l}\leq\psi_l+c_l$. Choose
$0<\delta<\min_l\delta_l$ satisfying the new target gap. For
$b=(1-\delta)\psi-k$, at every finite old value,
\[
 b-F_l\leq(\delta_l-\delta)\psi_l+(1-\delta)c_l+k_l-k.
\]
One large $k$ makes $b\leq F_l$ for all $l$; at an old pole the same
reference bound forces $b=-\infty$. Choose
$0<\kappa<\min_l\kappa_l/4$. Apply the compact comparison extension
first to $u_i+2\kappa$, then to a late defining tail of this fixed $u_i$,
using (4.2). One obtains an actual stage $h$ sufficiently close to $L$
to tensor into new $T_\eta$, and $h+\kappa\leq F_l$ for all $l$. Then
\[
 F=\max(h+\kappa,b)\leq F_l\quad(\forall l),\qquad F\geq u+\kappa.
\]
All special stages, ordered by scalar decrease and by $\eta$ decrease,
form a directed net. Its class tends to $\theta$. Its infimum is $u$
at every point: for $t>u(x)$ finite choose $u_i(x)<t$, select the target
gap before a sufficiently late defining stage $h$, make the guard small,
and make the floor smaller than $t$ at $x$. Old constraints can then be
imposed by directedness. This works also at $u(x)=-\infty$ for each finite
$t$. Thus the endpoint has one literal defining net. This proves genuine
decreasing closure with a finite generic endpoint.

For an upper bounded usc obstacle with a genuine generically finite
minorant, continuous majorants have pointwise infimum that obstacle on
compact Hausdorff $K$. Continuous cutoffs and a global upper bound prove
this even at $-\infty$. Uniform rational-model approximation and finite
minima give a directed family of rational model majorants. Their genuine
envelopes from (3.1) decrease, with a finite generic lower bound. Closure
makes their endpoint genuine; the obstacle inequality and every competitor
show it is exactly the envelope. This also gives cofinal decreasing
model-envelope majorants for every genuine function.

For bounded continuous $f$, Lemma \ref{sections} now gives
\[
 P_\theta(f)=Q_L(f)\quad\text{on }K\setminus Z.
\tag{4.3}
\]
The right side is lsc there as a supremum of extended-continuous FS
functions; the left side is usc by genuineness. Both are finite there.
Hence every such envelope is continuous on this ordinary open set.
For model obstacles the full relative-system stages increase on a
cofinal factorial sequence and their supremum is $Q_L(f_D)$ by G1.
Both discrepancies $P_\theta(f_D)-Q_L(f_D)$ and $R-Q_L(0)$ vanish off
the same $Z$. No assertion about their values on $Z$ is needed below.

\section{Actual scalar energy and its full finite-energy space}
Here is the scalar construction used in the target theorem. Represent a
rational model obstacle $f_D$ on a smooth equivariant modification
$Y\to X\times\mathbb P^1$, with full central fiber $F=\sum b_EE$ and
vertical $D$, so $f_D(v_E)=d_E/b_E$. Choose $aF\leq D\leq bF$.
Proper normal birational section descent identifies its relative full
section lattice, with $R_m=H^0(X,mL)$, as
\[
 t^{-ma}R_m[t]\subset\Lambda_m\subset t^{-mb}R_m[t],\quad
 \Lambda_m=\bigoplus_j t^{-w_{mj}}k[t]s_{mj},\quad ma\leq w_{mj}\leq mb.
\]
The last equality follows by a basis adapted to the finite homogeneous
flag, not finite generation of the whole section ring. Gluing to the
product lattice at infinity gives
\[
 h^0(Y,m(\rho_X^*L+D+cF))=\sum_j\max(w_{mj}+mc+1,0).
\]
For fixed rational $c>-a$, actual section-volume limits give
\[
 \vol_Y(\rho_X^*L+D+cF)=(n+1)cV+(n+1)\mathscr E(f_D).
\tag{5.1}
\]
The intercept is the limit of $n!m^{-n-1}\sum_jw_{mj}$. Refining models
preserves the actual lattice. It is monotone, $V$-translation equivariant,
and $V$-Lipschitz for the original sup norm. Volume-root concavity on
fixed big algebraic classes, expanded at $c\to\infty$ in (5.1), gives
concavity of this intercept on rational model functions. Uniform model
density gives its unique concave monotone $V$-Lipschitz extension to $C(K)$.

This extension is also the actual full-section trace. For $0\ne s\in R_m$
put $W_m(f;s)=\inf_K(mf-\log\|s\|)$; at section zeros the bracket is
$+\infty$. Its values lie in $[m\min f,m f(v_0)]$. Upper-weight spaces
are linear by the ultrametric inequality and their endpoints are closed.
A finite flag therefore gives an adapted basis with trace
\[
 e_m(f)=n!m^{-n-1}\sum_j\lfloor W_m(f;s_{mj})\rfloor.
\tag{5.2}
\]
For model obstacles these floors are exactly the lattice weights above.
For arbitrary continuous $f,g$ flag comparison gives
$|e_m(f)-e_m(g)|\leq(n!\dim R_m/m^n)(\|f-g\|_\infty+1/m)$.
Thus (5.2) converges to $\mathscr E(f)$ in cofinal admissible degrees.

Define on the entire actual cone
\[
 E(u)=\inf_{f\in C(K),\ f\geq u}\mathscr E(f),\qquad
 \E^1(\theta)=\{u\in\PSH(\theta):E(u)>-\infty\}.
\tag{5.3}
\]
For every section $s$, a weighted section scalar lies below $f$ iff it
lies below $P_\theta(f)$. Hence its full weights for these two obstacles
are identical. If $g\in C(K)$ dominates $P_\theta(f)$, weight monotonicity
gives $\mathscr E(g)\geq\mathscr E(f)$. Taking the outer infimum proves
$E(P_\theta(f))=\mathscr E(f)$. In particular
\[
 E(R)=0,\quad E(u+c)=E(u)+cV,
 \quad B(u):=-E(u)\geq0\text{ for }u(v_0)=0.
\]
These formulas use section volume, not $L^n$ in a nonnef class.

Extend (5.3) to all upper bounded usc obstacles. If $u_i\downarrow u$
and $f\geq u$ is continuous, compactness of
$\{u_i\geq f+\epsilon\}$ gives a tail $u_i<f+\epsilon$. Monotonicity
then gives $\inf_i E(u_i)\leq\mathscr E(f)+\epsilon V$.
Taking the outer infimum proves decreasing-net energy continuity,
including $-\infty$. Continuous majorant nets of two usc obstacles,
their sums, and continuous concavity give concavity of the outer extension
and hence of $E$ on the genuine cone. The genuine floors
$\max(u,R-N)$ decrease all-point to normalized $u$ and have costs at
most $NV$. Their decreasing energies characterize all of (5.3); no bounded
relative-singularity restriction replaces this space.

This section constructs scalar energy, not a mixed unbounded operator.
In particular a polar physical complete-system function is not put into
$\E^1$ by its moving decomposition. Even for $\OO(1)$ on $\mathbb P^1$,
the floors of $\log\|z_1\|$ have energy $-N/2$, so its energy is
$-\infty$ by G0. The common horizontal factor used later is only a
literal section identity.

\section{Capped divisors, a genuine finite-energy pole, and targets}
Let $D\geq0$ be an actual integral Cartier divisor and
$a_D=-\log\|s_D\|\in[0,+\infty]$. Its caps are ordinary continuous.
For rational $t>0$ with $L-tD$ big, and integer $N\geq1$, set
\[
 f_{t,N}=-\min(t a_D,N),\qquad u_{t,N}=P_\theta(f_{t,N}).
\]
The minorant $R-N$ ensures genuine existence. In denominator-cleared
degrees the actual injection
\[
 s_D^{mt}:H^0(X,m(L-tD))\hookrightarrow H^0(X,mL)
\]
has every nonzero image vector of full obstacle weight zero. Indeed
$\log\|s' s_D^{mt}\|\leq-mt a_D\leq m f_{t,N}$ everywhere, with literal
$-\infty\leq-mN$ at supported zeros, and equality at $v_0$. All other
weights lie in $[-mN,0]$. An adapted basis retains at least
$h^0(m(L-tD))$ zero weights. Flooring and section-volume limits prove
\[
 0\leq B(u_{t,N})\leq N[\vol(L)-\vol(L-tD)].
\tag{6.1}
\]
The same image section is a genuine minorant of generic value zero, so
$u_{t,N}$ is normalized without a shift altering its cap. It satisfies
$R-N\leq u_{t,N}\leq R$ and $u_{t,N}\leq-N$ on $D^{\rm an}$.
The cap can also be represented as a rational model function: for
$t=a/b$, use $\OO(-aD)+(\tau^{bN})$ on $X\times\mathbb A^1$ and scale
its ideal model function by $1/b$. No approximation of supported poles
by finite valuations changes the calculation.

Choose $\lambda_j=2^{-j}$, $N_j=4^j$, and rational $t_j$ small enough that
$L-t_jD$ is big and $N_j[V-\vol(L-t_jD)]\leq2^{-j}$.
Volume continuity is used inside the big cone, with $N_j$ fixed first.
The actual convex functions
\[
 w_q=\sum_{j\leq q}\lambda_j u_{t_j,N_j}+2^{-q}R
\]
are normalized and decrease all-point. Concavity gives $E(w_q)\geq-1/3$.
The closure and energy continuity proved above give a genuine
$w_D\in\E^1(\theta)$ with $E(w_D)\geq-1/3$. On the entire original
$D^{\rm an}$ it is $-\infty$, since
$w_q\leq-\sum_{j\leq q}2^j$. This verifies the convex-series and
normalization steps proposed in Root197.

An intrinsic target here is a positive Radon measure of raw mass $V$ with
\[
 \nu\{R=-\infty\}=0,\qquad
 I_\nu(u):=\int(R-u)\,d\nu<\infty
 \quad\text{for every normalized actual }u\in\E^1(\theta).
\tag{6.2}
\]
The difference uses actual values where $R$ is finite. Applying (6.2) to
$w_D$ gives $\nu(D^{\rm an}\setminus\{R=-\infty\})=0$; the first
condition handles the rest. Every proper algebraic closed subset lies
in an actual divisor: twist its ideal by a sufficiently large actual
ample power and take a nonzero section. Hence every target (6.2) is
algebraically null. No absolute integrability of $R$ was used.
If one specifies finite intrinsic dual energy
\[
 D_\nu=\sup_{u(v_0)=0}\{E(u)+I_\nu(u)\}<\infty,
\tag{6.3}
\]
it directly implies (6.2) because every finite-energy cost is finite.
Conversely proving full equivalence and coercivity would use a
relative-contraction estimate; such a transfer is not an inverse equation.
The nullness conclusion alone does not concern all nonalgebraic polar or
negligible sets, nor continuity on varying weakly convergent tests.

\section{Corrected finite-nef identities and an ordered energy-gap bound}
On any fixed actual nef generic rational bundle $A$ let $T_\xi=(A,\xi)$
for a bounded relatively nef model scalar. Finite mixed model products
are positive, of raw mass $V_A=A^n$, and keep coefficients
$b_E(T|_E)^n$. Define $e_A(\xi)=T_\xi^{n+1}/(n+1)$ and
$J_A(\xi)=\xi(v_0)V_A-e_A(\xi)$, and
\[
 \|f\|_{T_1\cdots T_{n-1}}^2=-(0,f)^2T_1\cdots T_{n-1}\geq0,
 \quad d_A(\xi,\eta)=\max_{j<n}\|\xi-\eta\|_{T_\xi^jT_\eta^{n-1-j}}^2.
\]
Vertical Hodge and Cauchy--Schwarz are finite algebraic facts; adding
actual ample slots and taking their finite polynomial limits is another
justification through G0. No unbounded nonample estimate is invoked.

For $S=(0,\xi-\eta)$, two telescoping sums give exactly
\[
 T_\eta^{n+1}-T_\xi^{n+1}+(n+1)S T_\eta^n
 =-\sum_{j=0}^{n-1}(n-j)S^2T_\xi^jT_\eta^{n-1-j}.
\tag{7.1}
\]
For each $j$ the summand occurs for $k=j+1,\ldots,n$. In particular
\[
 J_A(\xi)=\frac1{n+1}\sum_{j<n}(n-j)\|\xi\|_{T_\xi^jT_0^{n-1-j}}^2.
\tag{7.2}
\]
The order $j+1$ with these factors is false. On
$\operatorname{Bl}_{(p,0)}(\mathbb P^2\times\mathbb P^1)$ take
$T_\xi=H-aE$, $0<a<1$. The actual relatively ample model has
$H^3=H^2E=HE^2=0$, $E^3=1$. Thus
$J=a^3/3$, $\|\xi\|_{T_0}^2=0$, $\|\xi\|_{T_\xi}^2=a^3$;
the erroneous order gives $2a^3/3$. This is a proof-level intersection
witness; the accompanying finite arithmetic only checks transcription.

\begin{lemma}[Ordered finite-model gap]\label{gap}
If $\xi\geq\eta$ on the fixed carrier, then
\[
 d_A(\xi,\eta)\leq n(n+1)[e_A(\xi)-e_A(\eta)].
\tag{7.3}
\]
\end{lemma}
\begin{proof}
Write $\Delta=e_A(\xi)-e_A(\eta)$. The energy difference is
\[
 \Delta=\frac1{n+1}\sum_{k=0}^n
 \int(\xi-\eta)\,T_\xi^kT_\eta^{n-k}.
\]
Every summand is nonnegative and each endpoint integral is at most
$(n+1)\Delta$. The tangent remainder in (7.1), divided by $n+1$, is
$-\Delta+\int(\xi-\eta)d\MA_A(\eta)\leq n\Delta$.
Its coefficient of each seminorm is at least $1/(n+1)$.
This proves (7.3). No generic normalization or division by $V_A$ is needed.
\end{proof}

For completeness the corrected finite slot-change argument gives
\[
 \|\xi-\xi'\|_{T_\psi^jT_0^{n-1-j}}^2
 \leq C_n d_A(\xi,\xi')^{2^{1-n}}
       \max\{J_A(\xi),J_A(\xi'),J_A(\psi)\}^{1-2^{1-n}}.
\tag{7.4}
\]
Here is the proof structure needed to fix its scope. Average slots show
$d_A(\xi,\xi')\asymp_n\|\xi-\xi'\|_{T_{(\xi+\xi')/2}^{n-1}}^2$.
The tangent remainder in (7.1) is comparable to $d_A$ and convex in its
first potential. Put $f=\xi-\xi'$, $\tau=(\xi+\xi')/2$,
$a=d_A(\xi,\xi')$, $b=\max(d_A(\xi,\eta),d_A(\xi',\eta))$ and
$b_j=\|f\|_{T_\eta^jT_\tau^{n-1-j}}^2$. Then $b_0\leq C_na$ and
$b_{n-1}\leq4b$. Replacing one slot, splitting the second $f$ into
$T_\xi-T_{\xi'}$, and using Cauchy--Schwarz with
$T_\xi,T_{\xi'}\leq2T_\tau$ gives
$b_{j+1}-b_j\leq C_n\sqrt{b_j b}$. If $a\leq b$, induction gives
$b_j\leq C_n a^{2^{-j}}b^{1-2^{-j}}$; if $a>b$ use $b_{n-1}\leq4b$.
Apply this bound through a third potential to derive
$d\leq C_n b^{2^{1-n}}\max(d,b)^{1-2^{1-n}}$; solving it proves
the quasi-triangle bound. Average arbitrary nef slots, expand the
nonnegative multinomial terms, use this bound, and (7.2); this proves
(7.4). For $n=1$ the empty-slot seminorm gives the result directly.

With $\alpha=2^{-n}$, telescoping the $n$th powers from $j=0$ and
applying (7.4) and Cauchy--Schwarz yields
\[
 \left|\int(\xi-\eta)d\MA_A(\psi)
       -(\xi(v_0)-\eta(v_0))V_A\right|
 \leq C_n d_A(\xi,\eta)^\alpha J_A(\psi)^{1/2}
       M^{1/2-\alpha},
\tag{7.5}
\]
where $M=\max(J_A(\xi),J_A(\eta),J_A(\psi))$.
Indeed each $\psi$ seminorm in that sum is at most
$d_A(\psi,0)^{1/2}$, and each difference seminorm is bounded by
the square root of (7.4). Equations (7.3)--(7.5) are the proposed
ordered-gap refinement, rather than a bound only in terms of the test's
total energy. They hold also when $A^n=0$: all raw positive top measures
then vanish, and no probability normalization appears.

\section{The explicit conditional output package F}
The following additional assertions are hypotheses of Sections 9--10.
Their presence does not turn them into caller grants. They are the part
of the unreviewed dependence chain not independently completed here.
\begin{enumerate}
\item[F1.] A target-independent map $M(u)$ is constructed on all actual
$\E^1(\theta)$, with positive original Radon outputs of mass $V$ and
constant invariance. For decreasing normalized model-envelope majorants
$u_i=P_\theta(f_i)\downarrow u$, where $f_i\leq0$ are decreasing rational
model functions and $f_i(v_0)=0$, $E(u_i)\downarrow E(u)$ and
$M(u_i)\to M(u)$ weakly. Each model output is the finite divisorial
measure $\mu_D=\sum b_E\langle(\rho_X^*L+D+cF)^n\rangle\cdot E\,
\delta_{v_E}$ for sufficiently large $c$. It is the derivative of
the scalar model energy. The normalized cone has the compactness claimed
in the intended foundations. That last assertion is not needed for the
tail theorem itself.
\item[F2.] Resolving the full $|mL|$ system and each of any fixed finite
list of complete relative systems gives, on common actual models,
$p_m^*L=A_m+T_m$, $mA_m$ generated, $T_m\geq0$, and bounded relatively
nef $\psi_{m,D}$ with
$\phi_{m,D}=q_{T_m}+\psi_{m,D}$ and $\phi_{m,0}=q_{T_m}$.
The full pushed raw model measures converge to $\mu_D$;
$A_m^n\to V$, $e_{A_m}(\psi_{m,D})\to E(P_\theta(f_D))$,
and $\psi_{m,D}(v_0)\to0$ for normalized obstacles.
For ordered $f_H\geq f_J$ the literal full-lattice inclusion gives
$\psi_{m,H}\geq\psi_{m,J}$. Full mixed moving products converge
against every fixed signed Cartier test, with its nef decomposition fixed
before degree. No multilinearity of ordinary big positive products is assumed.
\item[F3.] For nonpositive model obstacles of cost at most $B$, the finite
time-deficient derivative $\mu_{D,c}$ of
$g_D(c)=\vol(\rho_X^*L+D+cF)/(n+1)$ is positive, and
$0\leq\mu_{D,c}\leq\mu_D$ coefficientwise with raw mass deficit
at most $B/c$, whenever $c>B/V$. Original constant and continuous
directions differentiate the outer energy at every genuine $\E^1$
base, with both signs and output $M(u)$. The proposed construction of
this package in Full151 uses ordinary big volume differentiation,
not the ample unbounded four-point theorem. That entire construction
still needs an independent complete audit.
\end{enumerate}
The actual common moving factors in F2 are not original $\E^1$ objects.
For energy transfers from a fixed moving subspace one must retain
$(\min f)(V-A_m^n)$ for a fixed original continuous majorant $f$;
first its section-degree limit, then $m\to\infty$, then the outer
infimum. F2 includes the resulting energy compatibility, not a dismissal
of the finite error. No hypothesis in F is an inverse equation or strong
fixed-target integration theorem.

\paragraph{Finite-model reconstruction still needed for F.}
Ordinary complete-moving cofinality can be proved by fixing rational
nef competitors $\beta_i$ on one determination, choosing one small
ample $Q_0$ with $N_i-Q_0$ big, and taking
$Q_i=(1-\epsilon)\beta_i+\epsilon Q_0$. Both $Q_i$ and $N_i-Q_i$
have the required ample/big classes. In a common cleared degree a
nonzero section of the latter injects the generated full $qQ_i$
system into $|qN_i|$, hence its moving part dominates $Q_i$.
Let $\epsilon\downarrow0$ on the fixed determination. BFJ's directed
compactness identifies the limit with the positive product. A fixed
signed test is decomposed into two fixed nef tests before degree.
This algebraic argument explains F2's moving product component, but
does not itself prove F1 or the unbounded output identifications.

\section{Conditional gap transfer, uniform collars, and all-test tails}
All assertions in this section assume F, in addition to the earlier
proved scalar foundations. We give the actual measure passage.

For fixed normalized model envelopes $a=P(f_H)\geq b=P(f_J)$ from
ordered obstacles, and normalized model output base $w=P(f_D)$, F2,
Lemma \ref{gap}, and (7.5) imply
\[
 \int(a-b)\,dM(w)
 \leq C_n[E(a)-E(b)]^\alpha B(w)^{1/2}
       \max\{B(a),B(b),B(w)\}^{1/2-\alpha}.
\tag{9.1}
\]
The generic difference error is retained at each finite moving stage and
vanishes only after $m\to\infty$. To justify identifying its left side,
one can subtract the two bounded relative section-quotient integrals.
Their endpoint envelopes and reference are continuous off $Z$ by (4.3).
The original output $M(w)$ is finite divisorial and misses $Z$.
Fixed-degree bounds
\[
 \max(-C,\phi_{q,H}-R)\leq\phi_{m,H}-\phi_{m,0}
 \leq\min(C,a-\phi_{q,0})
\]
are respectively lsc and usc near these atoms. They converge there to
$a-R$ by (4.3). A continuous cutoff equal to one on their finite support
and supported in the finite-stage open set allows Radon-net
semicontinuity to squeeze the full integral. The omitted measure mass
tends to zero by weak convergence and the cutoff. This proves the
bounded weighted identification without uniform quotient convergence
and without subtracting supported infinities. The same argument applies
to $b-R$. Thus (9.1) follows with no atomic support assumption on the
domain of the ultimate output $M(u)$.

For each actual effective divisor $S$ there is a uniform algebraic collar
estimate, stronger than pointwise algebraic nullness. Put
$a_S=-\log\|s_S\|$. On a fixed determination of $S$ and fixed ample
$T$, choose $d$ with $dT-\rho_X^*S$ nef. For $D\leq0$ and fixed $c>B/V$,
write $\gamma=\rho_X^*L+cF$. The vertical part of the pulled $S$ has
coefficients $b_Ev_E(s_S)$. Its horizontal remainder is effective;
the ordinary volume derivative in each effective direction is nonnegative.
Positive-product order paired with the fixed nef test gives
\[
 \int a_S\,d\mu_{D,c}
 \leq\langle(\rho_X^*L+D+cF)^n\rangle\cdot\rho_X^*S
 \leq d\langle\gamma^n\rangle\cdot T=:C_{c,S}.
\]
F3's positive tail implies, for the original continuous cap,
\[
 \int\min(a_S,N)\,d\mu_D\leq C_{c,S}+NB/c.
\tag{9.2}
\]
F1 passes (9.2) to every $M(u)$ with $B(u)\leq B$. Consequently
\[
 \sup_{B(u)\leq B}M(u)\{a_S\geq N\}
 \leq C_{c,S}/N+B/c.
\tag{9.3}
\]
First choose $c$, then $N$. Letting $N\to\infty$ and then
$c\to\infty$ also proves $M(u)(S^{\rm an})=0$. The fixed determination
and its nef test precede every model or measure limit. This argument
uses no signed positive-product monotonicity.

Every bounded continuous function on $K\setminus Z$, with bounded
Borel extension elsewhere, has converging integrals under a weak net
of these outputs converging to an algebraically null limit. To see it,
use a continuous cutoff $\chi_N$ equal to one for $a_{E_0}\leq N$ and
zero for $a_{E_0}\geq2N$. The product extends continuously by zero,
since it vanishes on a neighborhood of $Z$. For a fixed limiting measure
ordinary Radon regularity makes the collar error small; (9.3) makes it
uniform on bounded-cost output families. The continuous integral passes
through weak convergence with this fixed collar. This proves bounded
relative integral convergence for all the model envelopes in (9.1).
It follows that (9.1) holds with $w$ replaced by every normalized
genuine output base $u\in\E^1$, by F1 and a fixed test pair.

Let $v_i=P(f_i)\downarrow v$ be normalized nonpositive model-majorant
envelopes of any $v\in\E^1$. For fixed $i$ use (9.1) on $v_i\geq v_j$,
then $j\to\infty$ with the output fixed. Monotone convergence on the
algebraically null complement gives
\[
 \int(v_i-v)\,dM(u)\leq
 C_n[E(v_i)-E(v)]^\alpha B(u)^{1/2}
       \max\{B(v),B(u)\}^{1/2-\alpha}.
\tag{9.4}
\]
Every difference is defined where $R$ is finite. Taking $i$ to the
zero obstacle first also gives the usual all-$\E^1$ relative moment
bound, with $B(v)^\alpha$ instead of the energy gap. Formula (9.4)
is the new approximation-tail bound: it tends to zero uniformly over
all outputs $B(u)\leq B$, for each fixed finite-energy test $v$.
For $n=1$ its last exponent is zero; the expression is read as
$C_1[E(v_i)-E(v)]^{1/2}B(u)^{1/2}$, including zero cases by continuity.

Fix such a family and $\epsilon>0$. Choose $i$ so its (9.4) error is
small. The bounded nonnegative function $R-v_i$ is continuous off $Z$,
with norm at most $C_i$. Fix $c$, then $N$ using (9.3), so its omitted
collar cost is small for this $C_i$. The actual original function
\[
 b=\chi_N(R-v_i)\text{ off }Z,\qquad b=0\text{ near }Z,
\tag{9.5}
\]
belongs to $C(K)$, is nonnegative, and obeys
\[
 \sup_{B(u)\leq B}\int|(R-v)-b|\,dM(u)<\epsilon.
\tag{9.6}
\]
Indeed the absolute error is bounded by
$v_i-v+(1-\chi_N)(R-v_i)$. Thus this is an absolute $L^1$ assertion
for an original continuous approximant, with the test index and polar
collar fixed before any output index. This also proves uniform integrable
tails for the fixed test over the output family: if $C=\|b\|_\infty$,
then $(R-v)\mathbf1_{\{R-v>2C\}}\leq2|(R-v)-b|$.
A bounded first-moment bound alone would not have implied it.
For a single fixed intrinsic target
(6.2), dominated convergence and its algebraic nullness give a small
error for the same construction after making $i,N$ sufficiently large;
this does not make the assertion uniform on arbitrary varying tests
against that target.

\section{Conditional weighted continuity and entire supporting planes}
Continue to assume F. The test-net ordering issue can be resolved without
assuming a defining net eventually lies below a model majorant at old
poles. Let $a\geq b$ be any normalized actual $\E^1$ pair. Choose
nonpositive model obstacles $f_\lambda\downarrow a$ and
$g_\tau\downarrow b$ as in Section 4. The ordered envelopes
\[
 A_{\lambda,\tau}=P_\theta(\max(f_\lambda,g_\tau)),\qquad
 B_\tau=P_\theta(g_\tau)
\]
decrease respectively to $a,b$, retain generic value zero, and have
decreasing energies. Apply (9.1) to this ordered pair at a fixed finite
model output. Its two separate integrals converge at the finite support;
all actual finite-energy tests there are integrable by the earlier
moment bound. No monotonicity of their difference is used. This gives
the gap estimate for every ordered actual pair at every model output.
For a fixed actual pair, (9.6) gives convergence of its two relative
integrals along canonical model-output majorants of any $u$. Therefore
\[
 \int(a-b)\,dM(u)\leq
 C_n[E(a)-E(b)]^\alpha B(u)^{1/2}
       \max\{B(a),B(b),B(u)\}^{1/2-\alpha}.
\tag{10.1}
\]
The order $a\geq b$ is retained intrinsically throughout this extension.

Decreasing weak continuity on all genuine bases also follows from F3,
rather than being a new assumption. For $u_j\downarrow u\in\E^1$ and
$h\in C(K)$ put $H_j(t)=E(P_\theta(u_j+th))$.
The outer envelope-energy identity, proved by continuous majorants and
decreasing closure, makes these finite concave functions of $t$.
Envelopes and their energies decrease to $H(t)=E(P_\theta(u+th))$ for
each fixed $t$. Their derivatives at zero are the corresponding
$M(u_j)$ integrals by F3. For $t>0$, concavity gives
\[
 \frac{H_j(t)-H_j(0)}t\leq\int h\,dM(u_j)
 \leq\frac{H_j(0)-H_j(-t)}t.
\]
Pass to the net limit at fixed $t$, then $t\downarrow0$. Both derivative
signs for $H$ squeeze the integrals to $\int h\,dM(u)$.
Thus $M(u_j)\to M(u)$ weakly on the original carrier.

If also $v_j\downarrow v\in\E^1$ is any normalized genuine decreasing
test net, energy continuity and (10.1) give
$\int(v_j-v)dM(u_j)\to0$, since the base costs are bounded.
For the fixed limiting test, (9.6) gives its weighted weak convergence:
multiply the continuous approximation there by an arbitrary fixed
$h\in C(K)$, and bound its absolute error by $\|h\|_\infty$ times the
$L^1$ error. Combining these two steps proves
\[
 (R-v_j)M(u_j)\longrightarrow(R-v)M(u)
\tag{10.2}
\]
weakly for all such genuine decreasing nets. It proves conditional
independence of these actual original relative weighted limits. It does
not identify every simultaneous mixed-moving quotient, whose physical
section test need not be in original $\E^1$, with this operator.

Whole supporting planes for the proposed output map do follow under F
plus decreasing weak continuity for the canonical base majorants. At a
model obstacle $f_D$, the derivative of the concave scalar energy is its
finite $\mu_D$, so for every actual $v\in\E^1$ the outer-majorant
argument gives
\[
 E(v)\leq\mathscr E(f_D)+\int(v-f_D)d\mu_D.
\]
Apply this to $a=P(f_D)$: $E(a)=\mathscr E(f_D)$ and $a\leq f_D$, so
$\int(f_D-a)d\mu_D=0$. Thus $\mu_D$ has full contact and
$E(v)\leq E(a)+\int(v-a)dM(a)$ for every such $v$.
All integrals at this finite divisorial support are well defined; the
relative moment theorem makes the actual tests finite at positive atoms.

Let $a_i=P(f_i)\downarrow u$ be canonical majorants. The integral
$\int(v-u)dM(a_i)$ converges to $\int(v-u)dM(u)$ by (9.6) applied
separately to the two fixed tests $v,u$. Also
$\int(a_i-u)dM(a_i)\to0$ by (9.4), uniformly over their bounded costs.
Passing the finite plane gives
\[
 E(v)\leq E(u)+\int(v-u)dM(u),\qquad u,v\in\E^1(\theta).
\tag{10.3}
\]
Normalization shifts are restored using the raw constant law and mass $V$.
This conditional deduction shows that every solution of $M(u)=\nu$, if
it exists for (6.2), maximizes the intrinsic objective: (10.3) makes
$E(v)+I_\nu(v)\leq E(u)+I_\nu(u)$.
It neither proves that a maximizer exists nor separates different maximizers.
We have not identified $M$ with all required unbounded mixed products, and
do not call (10.3) an unconditional genuine equation theorem.

\section{What still prevents completion of original P16.5}
The full positive-volume target remains every actual big rational $L$
on every smooth projective original $X$ over all the retained fields.
The following assertions are unproved here.
\begin{enumerate}
\item The entire target-independent outer-output construction F1--F3,
including its original continuous-direction derivatives with both signs,
needs a complete independent
audit. Sections 2--6 do not construct an unbounded mixed product operator.
\item Identify that map with the genuine original-carrier operator,
all mixed measures and every admissible simultaneous unbounded weighted
approximation. Moving quotients that are not actual $\E^1$ functions
are still outstanding; (10.2) only handles genuine original tests and
outputs of the conditional map.
\item Prove strong integration against each fixed intrinsic target
(6.2) or (6.3) on genuine weak cost-bounded sets, or prove a genuine
strong-completeness replacement excluding every energy/integral jump.
Uniform approximation (9.6) concerns output families; it does not concern
an arbitrary prescribed target. A fixed target's integration of all
actual tests does not by itself make their varying integrals continuous.
\item Actually attain $\sup(E+I_\nu)$, then justify stationarity with
both original continuous test signs and identify the target equation.
Weak usc of $E$ and a compact cluster do not attain this sum: $I_\nu$
has the wrong one-sided semicontinuity for that argument.
\item Prove original Borel strict-winning locality and comparison.
The set $\{u<v\}$ for two usc functions need not be open; one cannot
pass its indicator through weak measures. No nonalgebraic/negligible-set
nullness theorem has been obtained by discarding the common algebraic
divisor alone.
\item Prove strict separation and uniqueness modulo constants of
every genuine operator solution. Equation (10.3) only puts any such
solution among the maximizers. Uniqueness of a maximizer without an
entire-plane theorem would be insufficient; even with the conditional
plane, strict separation in a nonnef big class remains unproved.
\end{enumerate}
None of these is supplied by deleting ``ample'' from G0, by using
its unbounded estimates on the numerical nonnef class, by a finite
atomic inverse construction, or by a historical ledger. Physical
obstacles with the same genuine envelope are not distinct genuine
solutions and do not refute their uniqueness.

\paragraph{A diagnostic for the remaining target jump.}
On $[0,1]$ put $w(x)=1/x$ for $x>0$ and $w(0)=+\infty$, and
$\nu_m=m^{-1}\delta_{1/m}+(1-m^{-1})\delta_1$. They converge weakly to
$\delta_1$; all first moments of $w$ are bounded, but they tend to two
while the limiting moment is one. For any fixed cap $N>1$, the omitted
tail is one for $m>N$. This is a measure-theoretic diagnostic, not an
NA finite-energy counterexample. Conversely finite sums of integrable
test functions need not control all varying test tails without an
energy-gap or target approximation theorem. Neither diagnostic proves
that the requested nonample equation theorem is false.

\section{The precise additional zero-volume equation range}
For an actual numerically trivial rational $L$, including nontorsion
$\operatorname{Pic}^0$ bundles, use the same varying-actual-bundle
single-net cone with numerical class zero. Its reference is the canonical
zero scalar. For rational $\epsilon>0$, $L+\epsilon A$ is ample and its
weight-zero generating scalar is identically zero. This one constant
scalar net defines zero; shifts define every real constant.

Conversely normalize $u\in\PSH(0)$ at $v_0$ to zero. Integer tensoring
makes $m u\in\PSH(0)$ for all $m$. Embed every one into a single actual
ample cone. G0 normalized compactness gives a finite divisorial lower
bound $-C_z$ independent of $m$. Thus $-C_z\leq m u(z)\leq0$ implies
$u(z)=0$ at every rational divisorial valuation. G0 all-point determination
gives $u=0$ on the entire original carrier. Hence $\PSH(0)=\mathbb R$.
Positive ample perturbations have raw mass $\epsilon^n A^n\to0$;
every continuous integral is bounded by this mass times its sup norm.
Their unique limit is zero, and their constant energies tend to zero.
Thus raw energy is zero, $\E^1(0)=\mathbb R$, and the raw operator on
constants is zero. The sole positive mass-zero target is literally the
zero measure; the equation is attained and unique modulo constants.
Comparison and both energy derivative signs are trivial. No division
by zero occurs. This is the full numerical-zero rational-bundle range,
not the full nonbig pseudoeffective cone.

On $\mathbb P^1\times\mathbb P^1$, the nonzero nef
$p_1^*\OO(1)$ has zero volume and both constant and nonconstant bounded
pullback generated-system scalars. Adding $\epsilon\OO(1,1)$ gives
positive raw top masses $2\epsilon(1+\epsilon)\to0$ for both scalars.
Their zero top equations agree although the potentials differ by more
than a constant. Thus the raw top zero equation cannot retain uniqueness
on that entire cone. A different zero-volume datum must be specified
before claiming a broader theorem.

\appendix
\section{Primary versions, exact units, and read scope}
The raw primary connector responses are preserved in \texttt{sources/}.
\texttt{extract\_sources.py} extracts eight complete pertinent units;
\texttt{sources/unit-index.json} records the original renderer intervals,
their maps to inclusive one-based local LF coordinates, and byte hashes.
There are no missing renderer lines in those units and no conflicting
overlap text. Renderer lines and PDF page labels are not downloaded PDF
LF coordinates. The default inherited-network PDF attempt failed with
``Operation not permitted''; no original PDF bytes or hash are asserted.

The primary units used or independently reconstructed are:
\begin{itemize}
\item Boucksom--Jonsson, Addendum, arXiv:2206.07183v2, posted
6 September 2023, PDF date 7 September 2023: renderer 128--183 contains
the whole pertinent Lemma 1.5 statement/proof and actual FS convention;
198--257 contains the entire Lemmas 1.6--1.7 lattice statement/proofs.
These confirm the bounded-obstacle conclusion and full-system
calculation. We do not import its general Theorem A or its earlier
nonample decreasing-closure call as a new grant.
\item BFJ, arXiv:math/0608260v2, posted 4 September 2007: renderer
147 explicitly states the characteristic-zero algebraically closed
field scope; 319--336 gives nef-product positivity and monotonicity
with their proofs; 349--441 gives positive-product construction,
directed compactness, continuity, nef agreement and the psef-remainder
competitor statement. Renderer 448--545 includes volume and concavity
statements and the full pertinent volume-differentiation argument,
with both signs. The HTML's generated 2026 date is not a new version.
Resolution, section asymptotics, Fujita, Hodge index and the referenced
Morse inequality remain ordinary published background; no recursively
complete reading or reproof of all their proof literature is claimed.
\item BJ, arXiv:1801.08229v4, PDF dated 24 March 2022:
renderer 2705--2789 is the finite Hodge and concavity unit, and
2858--3107 is the complete finite-distance, slot-replacement,
quasi-triangle and mixed-slot proof unit. The printed coefficient
orientation is corrected by (7.1), independently of its provenance.
The quasi-triangle proof is read with its two different differences
through the third potential. No actual ample unbounded Lemma 7.30
is applied to a nonnef class here.
\end{itemize}
Exact G0/G1 and every unreviewed supplied body and ledger are preserved
as whole marked units in \texttt{supplied/}; their literal inclusive
assignment-LF coordinates are in \texttt{supplied/index.json}.
No path listed inside a historical source ledger was opened. No other
campaign, agent, credential or account state was inspected or changed.
The cloud runtime skill was read for network guidance only; the user
prohibited credential inspection. Its environment-status tool was not
available. No supervised compute facility was exposed and none was
needed; no resource-heavy computation was launched.

\section{Reproducibility and canonical result}
\texttt{frozen/input-original.json} preserves the supplied JSON bytes;
the decoded exact task and assignment are separately frozen. Run
\texttt{python3 rerun.py} in this directory. It verifies frozen input,
whole supplied units, all source maps, result support order, symbolic
coefficient identities and finite rational witnesses, then compiles
this paper twice with the existing pdfLaTeX. Exact commands, outputs,
environment versions and artifact hashes are retained in \texttt{logs/}
and the manifest. These are reproducibility and transcription checks,
not mathematical certification. The authoritative local notes include
complete proofs or explicit antecedents for their claims. Their support
entries name only earlier local notes.

Original P16.5 remains unsolved by this submission. The complete family,
original scope and completion criteria remain unchanged in the frozen
input. The canonical result is \texttt{candidate=false}.
\begin{thebibliography}{9}
\bibitem{ADD} S. Boucksom and M. Jonsson,
\emph{Addendum to global pluripotential theory over a trivially valued field},
\href{https://arxiv.org/pdf/2206.07183v2}{arXiv:2206.07183v2}.
\bibitem{BJ} S. Boucksom and M. Jonsson,
\emph{Global pluripotential theory over a trivially valued field},
\href{https://arxiv.org/pdf/1801.08229v4}{arXiv:1801.08229v4}.
\bibitem{BFJ} S. Boucksom, C. Favre and M. Jonsson,
\emph{Differentiability of volumes of divisors and a problem of Teissier},
\href{https://arxiv.org/html/math/0608260v2}{arXiv:math/0608260v2}.
\end{thebibliography}
\end{document}
